%! AP Statistics — Unit 9, Lesson 9.6 — Lesson Plan
\documentclass[10pt]{article}
\usepackage{apstats-article}
\usepackage{apstats-boxes}
\usepackage{graphicx}
\graphicspath{{images/}}
\newcommand{\UnitNumberName}{Unit 9: Inference for Quantitative Data: Slopes \quad}
\newcommand{\LessonNumberName}{Lesson 9.6: Skills Focus --- Selecting an Appropriate Inference Procedure}

\begin{document}
\begin{center}
    {\Large\bfseries \CourseName: \SchoolYear \\
    {\small\bfseries \UnitNumberName \LessonNumberName}}
\end{center}

\begin{tcolorbox}[colback=sky, colframe=navy, boxrule=0.9pt, arc=2mm,
    left=2mm, right=2mm, top=1.5mm, bottom=1.5mm]
    \textbf{Primary Objective:} Students will synthesize all inference procedures studied
    in Units 6--9 (proportions, means, chi-square, and slope) to correctly identify the
    appropriate procedure for any given context, distinguishing between confidence intervals
    and significance tests and recognizing the data structure that triggers each procedure.
    (Big Idea: UNC, VAR, DAT)
\end{tcolorbox}

\renewcommand{\arraystretch}{1.6}
\begin{skillbox}[Priority Ideas \& Skills]{goldbox}
\begin{minipage}[t]{0.48\linewidth}
    \textbf{AP Skills Addressed}
    \begin{itemize}
        \item Identify the appropriate inference procedure for a given context (Skill 1.D, 1.E).
        \item Distinguish confidence intervals from significance tests (Skill 1.D).
        \item Recognize categorical vs.\ quantitative variables and one-sample vs.\ two-sample
              vs.\ paired vs.\ bivariate data structures (Skill 1.A).
        \item Justify the choice of procedure with a well-reasoned explanation (Skill 4.B, 4.D, 4.E).
    \end{itemize}
\end{minipage}
\hfill
\begin{minipage}[t]{0.48\linewidth}
    \textbf{Key Understandings}
    \begin{itemize}
        \item The type of variable (categorical vs.\ quantitative) is the first decision point.
        \item The data structure (one sample, two independent samples, matched pairs, or bivariate)
              determines the specific procedure.
        \item CI $\to$ estimate a parameter; significance test $\to$ assess evidence against $H_0$.
        \item Chi-square procedures are always significance tests; slope inference applies to
              bivariate quantitative data.
        \item Correctly naming the procedure is an AP exam expectation.
    \end{itemize}
\end{minipage}
\end{skillbox}

\begin{skillbox}[{All Inference Procedures: Quick Reference}]{greenbox}
\renewcommand{\arraystretch}{1.8}
\begin{tabularx}{\linewidth}{@{} p{0.22\linewidth} p{0.22\linewidth} X X @{}}
  \rowcolor{sky}
  \textbf{Variable Type} & \textbf{Data Structure} & \textbf{Confidence Interval} & \textbf{Significance Test} \\
  \hline
  Categorical & 1 proportion & 1-prop $z$-interval & 1-prop $z$-test \\
  Categorical & 2 proportions & 2-prop $z$-interval & 2-prop $z$-test \\
  Categorical & 1 var, $\ge 3$ categories vs.\ claimed dist. & --- & Chi-square GOF test \\
  Categorical & 1 var, 2+ independent groups & --- & Chi-square test for homogeneity \\
  Categorical & 2 categorical vars, 1 sample & --- & Chi-square test for independence \\
  Quantitative & 1 mean & 1-sample $t$-interval & 1-sample $t$-test \\
  Quantitative & 2 independent means & 2-sample $t$-interval & 2-sample $t$-test \\
  Quantitative & Matched pairs & Paired $t$-interval & Paired $t$-test \\
  Quantitative & Bivariate (slope) & $t$-interval for slope & $t$-test for slope \\
\end{tabularx}
\end{skillbox}

\begin{skillbox}[Activate Prior Knowledge \& Spiral Review]{sky}
\begin{minipage}[t]{0.48\linewidth}
    \textbf{Prior Knowledge Check (3--4 min)}
    \begin{itemize}
        \item Units 6--7: z-procedures for proportions; t-procedures for means;
              conditions (Large Counts, Normal/Large Sample, Random, 10\% rule).
        \item Unit 8: chi-square GOF, homogeneity, and independence.
        \item Unit 9: regression output, the $t$-test and $t$-interval for slope.
        \item The warm-up asks students to name the procedure from a brief context ---
              a synthesis across all units.
    \end{itemize}
\end{minipage}
\hfill
\begin{minipage}[t]{0.48\linewidth}
    \textbf{Connections}
    \begin{itemize}
        \item This lesson synthesizes Units 6--9; there is no new content.
        \item AP exam free-response questions frequently require students to identify
              the correct procedure before setting it up.
        \item Completing this lesson prepares students for the Unit 9 sample test
              and the AP exam.
    \end{itemize}
\end{minipage}
\end{skillbox}

\begin{skillbox}[Hook \& Lesson]{sky}
\begin{minipage}[t]{0.48\linewidth}
    \textbf{Hook: ``Same Answer --- Right Procedure?'' (4--5 min)}\\
    Display three scenarios simultaneously:
    \begin{enumerate}
        \item ``Is the mean commute time for bus riders less than 30 minutes?''
        \item ``Do urban and suburban commuters differ in mean commute time?''
        \item ``Is there an association between commute mode (bus/train/car) and
              commute satisfaction (high/low)?''
    \end{enumerate}
    Ask: ``Every scenario is about commuting. How are the procedures different?''
    Take three 30-second responses; record the correct procedure next to each.
\end{minipage}
\hfill
\begin{minipage}[t]{0.48\linewidth}
    \textbf{The Decision Framework (15 min)}
    \begin{enumerate}
        \item \textbf{Variable type:} Is the variable of interest categorical (proportions/counts)
              or quantitative (means/slope)?
        \item \textbf{Data structure:}
              \begin{itemize}
                  \item One sample $\to$ 1-prop $z$, 1-sample $t$, or chi-sq GOF.
                  \item Two independent groups $\to$ 2-prop $z$, 2-sample $t$, or homogeneity.
                  \item Matched pairs $\to$ paired $t$.
                  \item Two categorical variables, one sample $\to$ independence.
                  \item Bivariate quantitative $\to$ slope procedures.
              \end{itemize}
        \item \textbf{Goal:} Estimate (CI) or test a claim (test)?
        \item \textbf{Name the procedure exactly.}
    \end{enumerate}
\end{minipage}
\end{skillbox}

\begin{skillbox}[Explicit Instruction: Conditions by Procedure]{sky}
\renewcommand{\arraystretch}{1.9}
\begin{tabularx}{\linewidth}{@{} p{0.28\linewidth} X @{}}
  \rowcolor{sky}
  \textbf{Procedure} & \textbf{Key Conditions (beyond Random / 10\% rule)} \\
  \hline
  1-prop / 2-prop $z$ & Large Counts: $np_0 \ge 10$ and $n(1-p_0) \ge 10$ (or $n\hat p \ge 10$, $n(1-\hat p) \ge 10$) \\
  1-sample / 2-sample $t$ & Normal/Large Sample: population normal OR $n \ge 30$ \\
  Paired $t$ & Normal/Large Sample applied to the \emph{differences} \\
  Chi-square (all) & Expected counts $\ge 5$ in every cell \\
  $t$-test / $t$-interval for slope & LINER: Linear, Independent, Normal residuals, Equal variance, Random \\
\end{tabularx}
\end{skillbox}

\begin{skillbox}[Active Monitoring]{redbox}
\begin{minipage}[t]{0.48\linewidth}
    \textbf{Circulate and Check}
    \begin{itemize}
        \item Students use the three-step framework (type, structure, goal) before
              naming the procedure.
        \item Watch for students who default to $t$-tests for all quantitative data;
              prompt: ``Are these two independent groups or the same subjects measured twice?''
        \item Ensure students do not confuse chi-square independence with homogeneity.
        \item Students should name the procedure precisely (e.g., ``2-sample $t$-test for
              $\mu_1 - \mu_2$''), not vaguely (e.g., ``a $t$-test'').
    \end{itemize}
\end{minipage}
\hfill
\begin{minipage}[t]{0.48\linewidth}
    \textbf{Cold-Call Prompts}
    \begin{itemize}
        \item ``The variable is salary. Is that categorical or quantitative?''
        \item ``The researcher randomly selected men and women separately --- are these
              paired or independent samples?''
        \item ``You want to \emph{estimate} the slope. Is that a CI or a test?''
        \item ``There are only two categories --- can we use chi-square? Why or why not?''
        \item ``State the parameter of interest for this scenario.''
    \end{itemize}
\end{minipage}
\end{skillbox}

\begin{skillbox}[Group Work \& Differentiation]{redbox}
\begin{multicols}{3}
    \textbf{Tier R --- Remediate}
    \begin{itemize}
        \item Procedure selection table and decision flowchart provided.
        \item Identify the variable type and data structure only; procedure name
              is filled in from the table.
        \item Focus on categorical vs.\ quantitative and one vs.\ two groups.
    \end{itemize}
    \columnbreak
    \textbf{Tier A --- Approaching}
    \begin{itemize}
        \item Classify all 8 scenarios independently.
        \item State the parameter being estimated or tested for each.
        \item No reference table.
    \end{itemize}
    \columnbreak
    \textbf{Tier E --- Extension}
    \begin{itemize}
        \item Complete all Tier A work.
        \item For each scenario, also identify what additional information would be
              needed to carry out the procedure and what conditions to check.
        \item One bonus scenario is ambiguous; students must argue for their choice.
    \end{itemize}
\end{multicols}
\end{skillbox}

\begin{skillbox}[Individual Work \& Assessment]{redbox}
\begin{minipage}[t]{0.48\linewidth}
    \textbf{Exit Ticket (5 min)}\\
    Two brief scenarios; for each, students name the correct procedure and write
    a 2-sentence justification. Scored for correct identification and reasoning.
\end{minipage}
\hfill
\begin{minipage}[t]{0.48\linewidth}
    \textbf{AP-Style Check}\\
    \textit{Researchers randomly sample 40 runners before and after a training program
    and record each runner's finishing time. They want to test whether the program
    reduced finishing times. Which procedure?}
    \begin{enumerate}[label=(\Alph*)]
        \item 2-sample $t$-test for $\mu_1 - \mu_2$
        \item Paired $t$-test for $\mu_d$
        \item 1-sample $t$-test for $\mu$
        \item $t$-test for slope
    \end{enumerate}
    Answer: (B) --- same runners measured twice; differences are the data.
\end{minipage}
\end{skillbox}

\begin{skillbox}[Reinforcement \& Extension]{goldbox}
\begin{minipage}[t]{0.48\linewidth}
    \textbf{Homework Overview}
    \begin{itemize}
        \item Five scenarios spanning all procedure families; students identify the
              procedure, state $H_0$/$H_a$ (or the parameter for a CI), and list
              the key conditions to check.
        \item At least one scenario involves slope inference and one involves
              paired data --- both commonly misclassified on AP exams.
    \end{itemize}
\end{minipage}
\hfill
\begin{minipage}[t]{0.48\linewidth}
    \textbf{Spiral Preview}
    \begin{itemize}
        \item This lesson completes the inference synthesis for Units 6--9.
        \item Students are now prepared for the AP-style Unit 9 test, which will
              require selecting and carrying out inference procedures across all
              quantitative and categorical contexts.
        \item Encourage students to keep the procedure selection table for exam review.
    \end{itemize}
\end{minipage}
\end{skillbox}

\end{document}
